\documentclass[10pt,a4paper]{amsart}
\usepackage[margin=19mm]{geometry}
\usepackage{amsmath,amssymb,mathtools}
\usepackage[scr=rsfs,cal=euler]{mathalfa}
\usepackage{xfrac}
\usepackage[unicode,hidelinks,bookmarks=false]{hyperref}
\newcommand{\ie}{i.e.\ }
\newcommand{\cf}{cf.\ }
\newcommand{\resp}{respectively\ }
\newcommand{\Subsection}{\subsection}
\providecommand{\nobreakdash}{\nobreak\hbox{-}}
\setlength{\emergencystretch}{2em}
\multlinegap0pt
\newcommand{\deltanH}{\delta_{n}}
\usepackage{amssymb}
\usepackage{natbib}
\usepackage[normalem]{ulem}
\usepackage{mathtools}
\usepackage{xcolor}
\newtheorem{lem}{Lemma}
\newtheorem{claim}{Claim}
\usepackage{cleveref}
\crefname{lem}{Lemma}{Lemmas}
\crefname{claim}{Claim}{Claims}
\newcommand{\bG}{\mathbf{G}} 
\newcommand{\bH}{\mathbf{H}} 
\newcommand{\bX}{\mathbf{X}} 
\newcommand{\bY}{\mathbf{Y}} 
\newcommand{\bbH}{\mathbb{H}} 
\newcommand{\bbP}{\mathbb{P}} 
\newcommand{\bbQ}{\mathbb{Q}} 
\newcommand{\dtv}{\mathrm{d}_{\mathrm{TV}}}
\newcommand{\e}{{\mathbf E}}
\title{The statistical threshold for planted matchings and spanning trees}
\author{Louigi Addario-Berry, Omer Angel, G\'abor Lugosi, Mikl\'os Z. R\'acz, Tselil Schramm}
\date{Source: arXiv:2602.07669v1 (CC BY 4.0)}
\begin{document}
\maketitle

\noindent\emph{Source and licence.} The statistical threshold for planted matchings and spanning trees. Version arXiv:2602.07669v1 (CC BY 4.0), distributed by arXiv under the Creative Commons Attribution 4.0 International licence, http://creativecommons.org/licenses/by/4.0/, as recorded in the arXiv licence metadata for this exact version and shown on the abstract page https://arxiv.org/abs/2602.07669v1. Full licence notice: \texttt{licenses/arxiv-2602.07669-CC-BY-4.0.txt}.\par\smallskip

\noindent\textit{Editorial scope.} Counted unit: the article's lem lem (source0.tex lines 388--390). The article's complete original proof of it is delivered with it (source0.tex lines 392--500, one \texttt{proof} environment of 109 lines). No mathematics is written, repaired or re-derived: the statement and the proof are the article's own lines, in the article's own notation, and no retained line is substituted or re-flowed.  Retained uncounted as the source-local context the proof needs: lines 366--370.  Nothing was omitted from any retained span, so the package carries no omission record.  No retained line contains a citation, so the wrapper carries no bibliography and no bibliography entry.  No retained line contains a figure environment, an image-inclusion command or drawing code, so no image is delivered and the package declares no asset dependency.  Editorial formatting only, in addition: an AMS \texttt{amsart} preamble that restates the article's own declarations and macros used by the retained lines, namely the environments \texttt{lem}, \texttt{claim} on separate counters, together with the standard packages those lines need.  The retained environments print in this document as Lemma 1, Claim 1 in order of appearance, whereas the article prints the same items with the numbering of its own full document.  The complete native source of the article as distributed in the arXiv e-print for this exact version is bundled unchanged as \texttt{sources/194197/source0.tex}.
\begin{align}
\e_{\bG \sim \bbQ} L(\bG)^k
&= 
\left(\frac{1}{1-\deltanH}\right)^{k\binom{n}{2}}\left(\frac{q_k}{p}\right)^{k e_H} Z_k \cdot  \e_{\bH_1,\ldots,\bH_k \sim \bbH} \left[q_k^{-\mathrm{collisions}(\bH_1,\ldots,\bH_k)}\right].\label{eq:coll-simple}
\end{align}

\begin{lem}
If $\bbH$ is the uniform distribution over perfect matchings of $K_n$ and $p = \omega(n^{-1/2})$ then $\lim_{n \to \infty} \dtv(\bbP,\bbQ) = 0$.
\end{lem}

\begin{proof}
It suffices to show that $\chi^2(\bbP,\bbQ) = \e_{\bG \sim \bbQ} (L(\bG)^2 - 1) =o_n(1)$.
We begin by instantiating \cref{eq:coll-simple}:
\begin{equation}\label{eq:coll-simple-matching}
\e_{\bG \sim \bbQ} L(\bG)^2 =
\left(\frac{1}{1-\deltanH}\right)^{2\binom{n}{2}}\left(\frac{q_2}{p}\right)^{n} \cdot Z_2 \cdot \e_{\bH_1,\bH_2 \sim \bbH} \left[q_2^{-\mathrm{collisions}(\bH_1,\bH_2)}\right].
\end{equation}
We next obtain an estimate of the moment generating function of $\mathrm{collisions}(\bH_1,\bH_2)$.
\begin{claim}\label{claim:pois-mgf}
    If $s n \gg \frac{\log n}{\log\log n}$, then
    \[
    \e_{\bH_1,\bH_2 \sim \bbH}\left[s^{-\mathrm{collisions}(\bH_1,\bH_2)}\right] \le \left(1 + o_n(1)\right) \e_{\bY \sim \mathrm{Pois}(\frac{1}{2})}\left[s^{-\bY}\right].
    \]
\end{claim}
\begin{proof}
    
We argue that the random variable $\bX = \mathrm{collisions}(\bH_1,\bH_2)$ is close to Poisson with mean $\frac{1}{2}$, such that the above is well-approximated by the Poisson moment generating function.
The set of perfect matchings of $K_n$ is in bijection with unordered partitions of $n$ into pairs, of which there are $(n-1)!!$.
By inclusion-exclusion, when $\bH_1,\bH_2$ are chosen independently and uniformly at random from this set, for any $k \in [n/2]$,
\begin{align*} 
\Pr [\bX = k] 
&= \frac{\binom{n/2}{k}\sum_{\ell=0}^{n/2} (-1)^\ell \cdot\binom{n/2 - k}{\ell} \cdot (n-2k - 2\ell-1)!!}{(n-1)!!}.
\end{align*}
We obtain two upper bounds on this probability.
First, for any $k \in [n/2]$ we obtain a crude bound from the $\ell=0$ summand:
\[
\Pr[\bX = k]
\le \binom{n/2}{k}\frac{(n-2k-1)!!}{(n-1)!!} 
= \frac{2^{n/2}}{k!\binom{n}{n/2}} \cdot \frac{1}{2^{n/2-k}} \binom{n-2k}{n/2-k} 
= O(\sqrt{n}) \cdot \frac{1}{2^k k!},
\]
where the final inequality follows from the fact that $\binom{2m}{m} = \frac{1}{\sqrt{\pi m }}2^{2m}(1+O(\tfrac{1}{m}))$.

Now, we obtain a finer bound for $k\ll n$.
If $k \ll n$, then letting $k^+$ be the smallest even integer at least as large as $k$, we over-estimate the probability by stopping the inclusion-exclusion at the $k^+$ term:
\begin{align*}
\Pr[\bX = k]
&\le \frac{\binom{n/2}{k}\sum_{\ell=0}^{k^+} (-1)^\ell \cdot\binom{n/2 - k}{\ell} \cdot (n-2k - 2\ell-1)!!}{(n-1)!!}\\
&= \frac{1}{k!} \frac{1}{\binom{n}{n/2}} \sum_{\ell=0}^{k^+} (-1)^\ell \frac{1}{\ell!} 2^{k+\ell} \binom{n-2k-2\ell}{n/2-k-\ell}\\
&= \frac{1}{k!} \frac{\sqrt{\pi n/2}}{2^{n}}\left(1+O(\tfrac{1}{n})\right) \cdot \sum_{\ell=0}^{k^+} (-1)^\ell \frac{1}{\ell!} \frac{2^{n-k-\ell}}{\sqrt{\pi(n/2-k-\ell)}}\left(1+O(\tfrac{1}{n})\right)\\
&= \frac{1}{2^k k!} \sum_{\ell=0}^{k^+} \left(-\frac{1}{2}\right)^\ell \frac{1}{\ell!}\left(1+O(\tfrac{k}{n})\right)\\
&\leq \left(1 + O(\tfrac{k}{n})\right) \cdot e^{-1/2} \frac{1}{2^k k!}.
\end{align*}

Let $\bY \sim \mathrm{Pois}(\frac{1}{2})$, with $\Pr[\bY = k] = e^{-1/2} \frac{1}{2^k k!}$.
Let $m \ll n$ be the smallest integer such that $s m > 100 \log n / \log\log n$.
We manipulate the moment generating function of $\bX$, applying the coarse bound when $k > m$ and the finer bound when $k \le m$ and obtain
\begin{align*}
    \sum_{k=0}^\infty \Pr[\bX = k] \left(\frac{1}{s}\right)^k
    &\le (1+o_n(1))\sum_{k=0}^{m} \Pr[\bY = k] \left(\frac{1}{s}\right)^k + \sum_{k=m+1}^n O(\sqrt{n})\frac{1}{2^k k!} \left(\frac{1}{s}\right)^k\\
    &\le (1+o_n(1))\sum_{k=0}^{m} \Pr[\bY = k] \left(\frac{1}{s}\right)^k + \sum_{k=m+1}^n  O(\sqrt{n k})\left(\frac{e}{k s}\right)^k\\
    &= (1+o_n(1)) \e\left[s^{-\bY}\right] + o_n(1)
    = (1+o_n(1))\e\left[s^{-\bY}\right]. \qedhere
\end{align*}
\end{proof}


In the case of a planted perfect matching we have that $\deltanH = \frac{1}{n-1}$, 
and a simple calculation shows that 
if $p = \omega(n^{-1/2})$, 
then $q_2 = \omega(n^{-1/2})$.  
Hence, \Cref{claim:pois-mgf} applies with $s = q_{2}$.
We next recall that when $\bY \sim \mathrm{Pois}(\lambda)$, then $\log(\e{s^{-\bY}}) = \lambda\left(\frac{1}{s} - 1\right)$.
Plugging this into \cref{eq:coll-simple-matching} and taking logarithms, we have that 
\begin{align}
    \log \e_{\bG \sim \bbQ} L(\bG)^2
    &\le o_n(1) - 2\binom{n}{2}\log\left(1-\deltanH\right)  + n\log\left(\frac{q_2}{p}\right)+ \log Z_2 + \frac{1}{2q_2} - \frac{1}{2}. \label{eq:log}
    \end{align}
    Some algebraic simplification gives us
\begin{align}
    q r_2 &= (q-\deltanH)(1-\deltanH/q),\label{eq:simplif}\\
    Z_2 &= (1-q+qr_2)^{\binom{n}{2}} = \left(1-2\deltanH + \tfrac{\deltanH^2}{q}\right)^{\binom{n}{2}},\nonumber\\
    q_2 &= Z_2^{-1/\binom{n}{2}} \cdot q r_2, \nonumber\\
    \frac{q_2}{p} &= (1-\deltanH)\left(1-\tfrac{\deltanH}{q}\right) Z_2^{-1/\binom{n}{2}}.\nonumber
\end{align}
So returning to \Cref{eq:log}, 
\begin{align*}
\log \e_{\bG\sim\bbQ} L(\bG)^2
&\le o_n(1) + \left(n - 2\binom{n}{2}\right) \log \left(1-\deltanH\right) + n \log\left(1-\tfrac{\deltanH}{q}\right)  \\
&\qquad \qquad + \left(1-\tfrac{n}{\binom{n}{2}}\right)\log Z_2 + \frac{1}{2 q_2} - \frac{1}{2}\\
&= o_n(1) - n(n-2)\log (1-\deltanH) + n \log(1-\tfrac{\deltanH}{q})  \\
&\qquad \qquad + \tfrac{n(n-3)}{2} \log \left(1 - 2\deltanH + \tfrac{\deltanH^2}{q}\right) + \frac{1}{2q_2} - \frac{1}{2}.
\end{align*}
Using the Taylor expansion of the logarithm $\log(1-x) = -x - \frac{x^{2}}{2} + o(x^{2})$ as $x \to 0$, 
and recalling that $\deltanH = \frac{1}{n-1}$ and $q = \omega(n^{-1/2})$, 
we have that 
$\log(1 - \deltanH) = - \deltanH - \deltanH^{2} / 2 + o( n^{-2})$, 
that 
$\log(1 - \deltanH / q) = - \deltanH / q + o(n^{-1})$, 
and that 
$\log(1 - 2\deltanH + \deltanH^{2} / q) 
= - 2 \deltanH + \deltanH^{2} / q - 2 \deltanH^{2} + o(n^{-2})$. 
Plugging these expressions into the display above, we see that the display above is equal to 
\[
o_n(1) 
- n(n-2)(-\deltanH - \tfrac{1}{2}\deltanH^2)
+ n \left(-\tfrac{\deltanH}{q}\right)
+ \tfrac{n(n-3)}{2}\left( -2\deltanH + \tfrac{\deltanH^2}{q} - 2\deltanH^2\right)
 + \frac{1}{2q_{2}} - \frac{1}{2}.
\]
This simplifies to 
\[
o_{n}(1) + \frac{1}{2q_{2}} - \left(1 + \Theta \left( \tfrac{1}{n} \right) \right) \frac{1}{2q}.
\]
Finally, using the fact that $q_{2} = q \left( 1 + \Theta\left( 1/(qn) \right) \right)$ 
and that $q = \omega(n^{-1/2})$, 
we obtain that the whole expression is $o_{n}(1)$. 
Hence $\chi^2(\bbP, \bbQ) = \e_{\bG \sim \bbQ} ( L(\bG)^2 - 1 ) = o_n(1)$, concluding the proof.
\end{proof}

\end{document}
